% =====================================================================
%  P1.  An unconditional bound for a Mobius-weighted correlation sum
%       in fixed residue classes, with two consequences for the
%       Huang-Li reduction of Goldbach to Elliott-Halberstam.
%
%  Deployment version.  Written to deploy/SPEC.md.
%  Compile:  pdflatex P1-mobius-fixed-class.tex   (twice)
% =====================================================================
\documentclass[11pt,a4paper]{article}

\usepackage{amsmath}
\usepackage{amssymb}
\usepackage{amsthm}
\usepackage[margin=1in]{geometry}
\usepackage[colorlinks=true,linkcolor=blue,citecolor=blue]{hyperref}

\theoremstyle{plain}
\newtheorem{theorem}{Theorem}
\newtheorem{corollary}[theorem]{Corollary}
\newtheorem{proposition}[theorem]{Proposition}
\newtheorem{lemma}[theorem]{Lemma}
\newtheorem{conjecture}[theorem]{Conjecture}
\theoremstyle{definition}
\newtheorem{observation}[theorem]{Observation}
\newtheorem{measurement}[theorem]{Measurement}
\theoremstyle{remark}
\newtheorem{note}[theorem]{Note}

\renewcommand{\SS}{\mathfrak{S}}
\newcommand{\AAA}{\mathfrak{A}}
\newcommand{\Emu}{E_\mu}
\DeclareMathOperator{\rad}{rad}
\DeclareMathOperator{\lcm}{lcm}

\title{An unconditional bound for a M\"obius-weighted correlation sum
in fixed residue classes, with two consequences for the Huang--Li
reduction of Goldbach to Elliott--Halberstam}
\author{}
\date{}

\begin{document}
\maketitle

\begin{abstract}
Huang and Li \cite{HL} show that the binary Goldbach conjecture for
large even $N$ follows from $EH$ together with a M\"obius-twisted
variant $EH_\mu$ whose levels sum to more than $1$; by
Bombieri--Vinogradov their Corollary~1 reduces this to
$EH_\mu(N^{\theta'})$ alone for a single $\theta'>1/2$. Their proof
consumes $EH_\mu$ at exactly two places, the functionals they call
$E_3(\alpha)$ and $E_4(\alpha)$. We prove unconditionally
(Theorem~\ref{thm:A}) that the M\"obius-weighted correlation sum
underlying $E_4$ is $\ll_A N(\log N)^{-A}$ for every $A>0$, uniformly
in the truncation point; consequently (Corollary~\ref{cor:B}) the
$E_4$ consumption of $EH_\mu$ is unnecessary and the entire $EH_\mu$
demand of the Huang--Li argument collapses to the single scalar
$E_3(\alpha)$. We then show (Theorem~\ref{thm:C}) that the
corresponding statement for $E_3$ --- the same sum with the weight
$\log k$ --- is not merely hard but, by an unconditional identity,
\emph{equivalent} to Huang--Li's own equation~(22): it therefore
already yields binary Goldbach for large even $N$, and yields the
asymptotic $\tilde r(N)\sim\SS(N)N$ exactly when
$\sum_{n<N}\Lambda(n)\mu(N-n)=o(N)$. The root cause is
$\mu*\log=\Lambda$. We then locate the true strength of the surviving
demand: Goldbach needs only a one-sided bound on $E_3$
(Proposition~\ref{prop:onesided}), at a threshold of order $N$ for
almost all $N$, and hence only a constant-factor bound on exactly the
object $EH_\mu$ bounds (Proposition~\ref{prop:nolog}). Finally we
record a defect in the published paper: equation~(18) drops an
$n$-dependent constraint present in the definition of $S_2(\alpha)$;
we give the missing term and show it is harmless under the hypotheses
already assumed (Section~\ref{sec:delta}).

We state plainly what is not claimed. Theorem~\ref{thm:A} removes only
the part of the demand that carries no Goldbach content, and the net
progress toward the Goldbach conjecture is zero. We also state plainly
what is standard: the ingredients of Theorem~\ref{thm:A} are classical
and its mechanism --- a divisor switch that moves the work onto a
short variable, followed by Bombieri--Vinogradov --- is ordinary
practice. What is offered is the application, together with the
corrections it turns out to require.
\end{abstract}

% =====================================================================
\section{Introduction}

\subsection{The Huang--Li reduction}

Let $N$ be a large even integer, $\Lambda$ the von Mangoldt function,
$\mu$ the M\"obius function, and
\[
  \tilde r(N)\;=\;\sum_{n<N}\Lambda(n)\Lambda(N-n),
  \qquad
  \SS(N)\;=\;2\prod_{p>2}\Bigl(1-\tfrac{1}{(p-1)^2}\Bigr)
              \prod_{\substack{p\mid N\\ p>2}}\Bigl(1+\tfrac{1}{p-2}\Bigr).
\]
Write $EH_\mu(Q)$ for the assertion that for every $A>0$
\begin{equation}\label{eq:EHmu}
  \sum_{q\le Q}\ \max_{y<N}\ \max_{(a,q)=1}
  \Bigl|\sum_{\substack{n\le y\\ n\equiv a\ (q)}}\Lambda(n)\mu(N-n)
        -\frac{1}{\varphi(q)}\sum_{n\le y}\Lambda(n)\mu(N-n)\Bigr|
  \;\ll_A\;\frac{N}{(\log N)^A}.
\end{equation}
Huang and Li \cite{HL}, continuing Pan \cite{Pan}, prove that
$EH(N^{\theta}(\log N)^{2A+8})$ together with $EH_\mu(N^{1-\theta})$
implies
\begin{equation}\label{eq:HL22}
  \tilde r(N)\;\ge\;\SS(N)\bigl(1-\AAA(N)\bigr)N
    +O\!\left(\frac{N}{(\log N)^A}\right),
\end{equation}
and deduce (their Corollary~1) that, in view of Bombieri--Vinogradov,
binary Goldbach for all sufficiently large even integers follows from
$EH_\mu(N^{\theta'})$ alone for any single $\theta'>1/2$. Here
\begin{equation}\label{eq:AN}
  \AAA(N)\;=\;\prod_{q\,\nmid\,N}\Bigl(1-\frac{1}{q(q-1)}\Bigr)
\end{equation}
is their equation~(7). We write $\AAA$ rather than their $A$ because
$A$ is already the free exponent in \eqref{eq:EHmu}, and one symbol
must not carry two meanings.

This is the sharpest conditional frame we are aware of: one sentence
stands between classical technology and Goldbach. This paper asks what
that sentence costs on the side where it is consumed.

\subsection{Where $EH_\mu$ is consumed}

The hypothesis \eqref{eq:EHmu} enters the Huang--Li proof at exactly
two places. With $\alpha$ the Vaughan-type parameter of their \S3 and
$K=(N-1)/\alpha$, these are (their notation; note that the residue
class is the \emph{fixed} class $a\equiv N$):
\begin{align}
  E_3(\alpha) &= \sum_{\substack{k<K\\ (k,N)=1}}\mu(k)\log k
    \Bigl[\sum_{\substack{n<N\\ n\equiv N\ (k)}}\Lambda(n)\mu(N-n)
      -\frac{1}{\varphi(k)}\sum_{n<N}\Lambda(n)\mu(N-n)\Bigr],
      \label{eq:E3}\\
  E_4(\alpha) &= \sum_{\substack{k<K\\ (k,N)=1}}\mu(k)
    \Bigl[\sum_{\substack{n<N\\ n\equiv N\ (k)}}\Lambda(n)\mu(N-n)\log(N-n)
      -\frac{1}{\varphi(k)}\sum_{n<N}\Lambda(n)\mu(N-n)\log(N-n)\Bigr].
      \label{eq:E4}
\end{align}
In both cases Huang--Li discard the signs $\mu(k)$ by the triangle
inequality on the first line and then appeal to \eqref{eq:EHmu} (their
Lemma~4 for $E_4$). The two functionals differ in exactly one respect:
$E_3$ carries the weight $w_k=\log k$ and $E_4$ carries $w_k=1$ (the
factor $\log(N-n)$ inside $E_4$ is $k$-independent and is removed by
partial summation). This paper shows that this one difference decides
everything.

\subsection{Results}

For $(k,N)=1$ and $1\le t<N$ put
\begin{equation}\label{eq:Emu}
  \Emu(t;k):=\sum_{\substack{n\le t\\ n\equiv N\ (k)}}\Lambda(n)\mu(N-n)
             -\frac{1}{\varphi(k)}\sum_{n\le t}\Lambda(n)\mu(N-n),
  \qquad
  C(t):=\sum_{n\le t}\Lambda(n)\mu(N-n),
\end{equation}
and, for a weight $w$,
\[
  T_w(t)\;:=\;\sum_{\substack{k<K\\ (k,N)=1}}\mu(k)\,w(k)\,\Emu(t;k).
\]
Thus $E_4$ is obtained from $T_1$ and $E_3$ from $T_{\log}$.

\begin{theorem}\label{thm:A}
Fix $\theta'\in(1/2,1)$ and put $K=N^{\theta'}$. Then for every $A>0$,
\[
  \sup_{1\le t<N}\bigl|T_1(t)\bigr|
  \;=\;\sup_{1\le t<N}
  \Bigl|\sum_{\substack{k<K\\(k,N)=1}}\mu(k)\,\Emu(t;k)\Bigr|
  \;\ll_{A,\theta'}\;\frac{N}{(\log N)^{A}}.
\]
Moreover every ingredient of the proof except Bombieri--Vinogradov
contributes only $O\!\left(Ne^{-c\sqrt{\log N}}\right)$: the power of
$\log$ in the statement is imposed by Bombieri--Vinogradov alone.
\end{theorem}

Theorem~\ref{thm:A} is a statement about a \emph{signed} sum in a
\emph{fixed} residue class. It is not an instance of, and does not
imply, $EH_\mu(N^{\theta'})$, which asks for absolute values and a
maximum over residue classes. The mechanism is described in
Section~\ref{sec:mechanism}: after divisor switching one
\emph{completes} the divisor sum, and on the complementary range
$k\ge K$ the M\"obius factor on the long variable squares itself away
while the surviving M\"obius sits on a variable of size
$\le N^{1/2-\delta}$. The completion is not a formality --- without it
the cofactor is not short and the configuration is the one with no
known machine.

\begin{corollary}\label{cor:B}
In the Corollary-1 regime, $E_4(\alpha)\ll_A N(\log N)^{-A}$ for every
$A>0$, unconditionally. Consequently Huang--Li's estimate
$S_4(\alpha)=O(N(\log N)^{-A})$ requires no hypothesis, their Lemma~4
is not needed, and the entire $EH_\mu$ demand of their argument
collapses to the single scalar $E_3(\alpha)$ --- together with the
correction $\Delta$ of Section~\ref{sec:delta}, which is likewise
unconditional there.
\end{corollary}

\begin{theorem}\label{thm:C}
In the Corollary-1 regime one has the unconditional identity
\[
  E_3(\alpha)\;=\;\sum_{n<N}\Lambda(n)\Lambda(N-n)
     \;-\;\SS(N)\Bigl(N-\sum_{n<N}\Lambda(n)\mu(N-n)\Bigr)
     \;+\;O_A\!\left(\frac{N}{(\log N)^{A}}\right).
\]
In particular the bound $E_3(\alpha)\ll_A N(\log N)^{-A}$ --- the form
of $EH_\mu$ that the Huang--Li argument consumes, though not the
weakest that suffices; see Proposition~\ref{prop:onesided} --- is
\emph{equivalent} to Huang--Li's equation~(22). It therefore yields
binary Goldbach for large even $N$ through their Theorem~1; and,
granted it, the asymptotic $\tilde r(N)\sim\SS(N)N$ holds if and only
if $C(N)=o(N)$.
\end{theorem}

\begin{note}[three statements, not two]\label{note:threeway}
The identity says that the displayed difference \emph{is}
$E_3(\alpha)$, so the hypothesis $E_3\ll_A N(\log N)^{-A}$ and
Huang--Li's~(22) are the same assertion. What \cite{HL} record
immediately after their~(22) is a different equivalence, between
$\tilde r(N)\sim\SS(N)N$ and $C(N)=o(N)$. Chaining the two,
$E_3\ll_A N(\log N)^{-A}$ delivers binary Goldbach --- Huang--Li's
Theorem~1 supplies $|C(N)|\le\AAA(N)N(1+o(1))$ by the triangle
inequality, and $\AAA(N)<1$ --- but it does not on its own deliver the
asymptotic, which needs $C(N)=o(N)$ in addition. The three statements
must be kept apart.
\end{note}

Theorem~\ref{thm:C} is the more consequential half. It says that the
demand side is closed at the level of \emph{identities} rather than of
\emph{estimates}: no choice of $\theta'$, of truncation, or of
smoothing can evade it, because the root cause is the identity
$\mu*\log=\Lambda$ from which Huang--Li start. The weight $\log k$ is
the carrier of the Goldbach content, so every divisor switch must hand
it back.

Closure at the level of identities is not closure at the level of
\emph{strength}, and the next two propositions separate the two. Every
condition on $E_3$ unwinds, through Theorem~\ref{thm:C}, to a
condition on $\tilde r(N)$ and $C(N)$ --- that is the identity part,
and it is permanent. But Huang--Li consume $E_3$ two-sidedly and at a
saving of every power of $\log$, and Goldbach needs neither.

\begin{proposition}[the demand is one-sided, and its threshold is not
a log power]\label{prop:onesided}
In the Corollary-1 regime, binary Goldbach for all sufficiently large
even $N$ follows from
\begin{equation}\label{eq:onesided}
  E_3(\alpha)\;>\;-\,\SS(N)\bigl(1-\AAA(N)\bigr)N\,\bigl(1+o(1)\bigr),
\end{equation}
with $\AAA(N)$ as in \eqref{eq:AN}. This is strictly weaker than
$E_3\ll_A N(\log N)^{-A}$ in two independent ways: it constrains $E_3$
from one side only, and its threshold is $\SS(N)(1-\AAA(N))N$, which
is $\asymp N$ for almost all even $N$ and never smaller than
$c\,N/(\log N\log\log N)$ for an absolute constant $c>0$.
\end{proposition}

\begin{proposition}[the demand needs no saving in $\log N$]%
\label{prop:nolog}
Put
\begin{equation}\label{eq:Bsum}
  B(N)\;:=\;\sum_{\substack{k<K\\ (k,N)=1}}(\log k)\,\bigl|\Emu(N;k)\bigr|,
\end{equation}
the quantity Huang--Li reach by the triangle inequality before
appealing to $EH_\mu$. Then binary Goldbach for all sufficiently large
even $N$ follows from
\begin{equation}\label{eq:nolog}
  B(N)\;\le\;(1-\varepsilon)\,\SS(N)\bigl(1-\AAA(N)\bigr)N
\end{equation}
for a single fixed $\varepsilon>0$. Since $\SS(N)(1-\AAA(N))\asymp1$
for almost all even $N$, \eqref{eq:nolog} is a \emph{constant-factor}
bound on exactly the object $EH_\mu$ bounds --- absolute values kept,
and no saving in $\log N$ whatsoever --- where $EH_\mu(N^{\theta'})$
asks for $\ll_A N(\log N)^{-A}$ for every $A$.
\end{proposition}

\subsection{What is not claimed}\label{sec:notclaimed}

\begin{itemize}
\item \textbf{No progress toward Goldbach.} Theorem~\ref{thm:A} is
unconditional but Goldbach-neutral: it removes exactly the half of the
demand that carries no Goldbach content. The net progress toward the
Goldbach conjecture is zero.

\item \textbf{No new estimate for $C(N)$.} Theorem~\ref{thm:C} places
the whole surviving demand on
$C(N)=\sum_{n<N}\Lambda(n)\mu(N-n)$; nothing here bounds that quantity.

\item \textbf{No claim of novelty of mechanism.} The ingredients of
Theorem~\ref{thm:A} are classical (Section~\ref{sec:lit}). Its
difficulty is bookkeeping, and the two places where the bookkeeping
bites are isolated in Section~\ref{sec:mechanism} and
Remark~\ref{note:trap}.

\item \textbf{The weakenings do not reopen the design space.}
Propositions~\ref{prop:onesided} and~\ref{prop:nolog} lower the
required strength but not the required level; the divisor-switch route
remains closed over its whole weight space \cite{P2}.

\item \textbf{The defect report is not a retraction.} \cite{HL}'s
Theorem~1 and Corollary~1 stand as stated; what
Section~\ref{sec:delta} supplies is a missing term and its treatment.
\end{itemize}

% =====================================================================
\section{Notation and the mechanism}\label{sec:mechanism}

$p$ and $q$ always denote primes. $\varphi$ is Euler's function,
$\tau_3$ the $3$-fold divisor function, $\rad(u)=\prod_{p\mid u}p$.
Constants $c,c_1,\dots$ are positive and absolute unless subscripted;
implied constants may depend on $A$ and $\theta'$. The letter $A$ is
reserved for the free exponent of \eqref{eq:EHmu}; the local factor of
\eqref{eq:AN} is $\AAA$.

The mechanism of Theorem~\ref{thm:A} is the following. The residue
class in \eqref{eq:Emu} is the fixed class $n\equiv N\pmod k$, i.e.
\[
  n\equiv N\pmod{k}\quad\Longleftrightarrow\quad k\mid N-n .
\]
So the $k$-sum is a divisor sum over $u=N-n$, namely the incomplete
sum $\sigma_K(u)=\sum_{k\mid u,\ k<K,\ (k,N)=1}\mu(k)$, and the
mechanism has three steps rather than one.

\emph{First} one \textbf{completes} it. By Lemma~\ref{lem:complete},
for squarefree $u$ the complete sum
$\sum_{k\mid u,\,(k,N)=1}\mu(k)$ equals $\mathbf 1_{\rad(u)\mid N}$,
which forces $u\mid\rad(N)$ and so leaves only $N^{o(1)}$ terms.
\emph{Then} the work sits entirely in the complementary sum, over
$k\ge K$ --- and it is only on that range that writing $u=mk$ gives
\[
  m\;=\;u/k\;<\;N/K\;=\;N^{1-\theta'}\;\le\;N^{1/2-\delta}.
\]
\emph{Finally}, for squarefree $u=mk$ the factorisation forces
$(m,k)=1$ and $\mu(u)\mu(k)=\mu(m)\mu^2(k)$: the M\"obius factor on
the \emph{long} variable $k$ cancels against itself and leaves the
nonnegative $\mu^2$, while the surviving M\"obius sits on the
\emph{short} variable $m$. That is the classical ``M\"obius on the
short variable plus Bombieri--Vinogradov'' configuration.

The completion step is load-bearing and easy to omit. Over the range
$k<K$ as it stands the cofactor $u/k$ runs up to $u$ itself, so the
bound $m<N^{1-\theta'}$ is simply false there and the surviving
M\"obius would sit on the \emph{long} variable --- the configuration
for which no machine is known. Steps~1--3 of Section~\ref{sec:proof}
carry out the argument in that order.

% =====================================================================
\section{Proof of Theorem~\ref{thm:A}}\label{sec:proof}

Fix $A>0$ and set
\begin{equation}\label{eq:trunc}
  D_0:=(\log N)^{A+2},\qquad E_0:=(\log N)^{2A+4}.
\end{equation}
The truncations must be allowed to depend on $A$; a fixed choice
covers only bounded $A$.

\subsection{Step 0: the subtracted mean term}

By \eqref{eq:Emu},
\begin{equation}\label{eq:split0}
  T_1(t)\;=\;D(t)\;-\;C(t)\,B(K),
  \qquad
  D(t):=\sum_{\substack{k<K\\(k,N)=1}}\mu(k)
        \sum_{\substack{n\le t\\ n\equiv N\ (k)}}\Lambda(n)\mu(N-n),
\end{equation}
where $B(K)=\sum_{k<K,(k,N)=1}\mu(k)/\varphi(k)$. Huang--Li's Lemma~1
(a form of Goldston--Y{\i}ld{\i}r{\i}m \cite{GY}) gives
$B(K)\ll e^{-c_1\sqrt{\log K}}$, and trivially
$|C(t)|\le\sum_{n<N}\Lambda(n)\ll N$. Hence
\begin{equation}\label{eq:meanterm}
  |C(t)B(K)|\;\ll\;Ne^{-c\sqrt{\log N}}
\end{equation}
uniformly in $t$.

\subsection{Step 1: divisor switching (an exact identity)}

Since $n\equiv N\pmod k$ with $n\le t<N$ is the same as $u:=N-n$ being
a multiple of $k$ in $[N-t,N)$,
\begin{equation}\label{eq:switch}
  D(t)\;=\;\sum_{N-t\le u<N}\Lambda(N-u)\,\mu(u)\,\sigma_K(u),
  \qquad
  \sigma_K(u):=\sum_{\substack{k\mid u,\ k<K\\ (k,N)=1}}\mu(k).
\end{equation}
This is a finite rearrangement, hence exact. It is verified
independently in
\texttt{audit\_switch\_identity.py}, which accumulates the two sides
in genuinely different index orders --- the $k$-side by modular
slices, the $u$-side against a divisor-sum array --- and finds them
equal to $10^{-16}$ relative to $N$.

\subsection{Step 2: the complete divisor sum}

\begin{lemma}\label{lem:complete}
For squarefree $u\ge1$,
$\displaystyle\sum_{k\mid u,\ (k,N)=1}\mu(k)=\mathbf 1_{\rad(u)\mid N}$.
\end{lemma}

\begin{proof}
The sum is multiplicative in $u$ with local factor $1$ at $p\mid N$
and $1-1=0$ at $p\nmid N$. Hence it vanishes unless every prime of $u$
divides $N$.
\end{proof}

Splitting $\sigma_K(u)$ into the complete sum minus the tail $k\ge K$,
\eqref{eq:switch} becomes $D(t)=P(t)-R(t)$ with
\begin{equation}\label{eq:PR}
  P(t)=\!\!\sum_{\substack{N-t\le u<N\\ \rad(u)\mid N}}\!\!\Lambda(N-u)\mu(u),
  \qquad
  R(t)=\!\!\sum_{N-t\le u<N}\!\!\Lambda(N-u)\mu(u)
       \sum_{\substack{k\mid u,\ k\ge K\\ (k,N)=1}}\mu(k).
\end{equation}
In $P(t)$ the conditions $\mu(u)\ne0$ and $\rad(u)\mid N$ force
$u\mid\rad(N)$, so there are at most $2^{\omega(N)}=N^{o(1)}$ terms,
each $\ll\log N$:
\begin{equation}\label{eq:P}
  P(t)\ll N^{o(1)}.
\end{equation}

\subsection{Step 3: the residual, and the degeneracy lemma}

Write $u=mk$ with $k\ge K$, so $m=u/k<N/K=:M$. For squarefree $u$ the
factorisation forces $(m,k)=1$ and $\mu(u)\mu(k)=\mu(m)\mu^2(k)$, so
\begin{equation}\label{eq:R1}
  R(t)=\sum_{m<M}\mu(m)
      \sum_{\substack{k\ge K,\ (k,m)=1,\ (k,N)=1\\ N-t\le mk<N}}
      \mu^2(k)\,\Lambda(N-mk).
\end{equation}

\begin{lemma}[degeneracy]\label{lem:degen}
The contribution to \eqref{eq:R1} of the terms with $(k,N)>1$ is
$O(N^{o(1)})$. The same bound holds for the terms in which the modulus
constructed in Step~4 is not coprime to $N$.
\end{lemma}

\begin{proof}
Suppose $p\mid(k,N)$ and $\Lambda(N-mk)\ne0$, say $N-mk=q^\ell$. Since
$p\mid k$ and $p\mid N$ we get $p\mid N-mk$, hence $q=p$ and
$n:=N-mk=p^\ell$ with $p\mid N$. The number of such $n<N$ is
$\le\sum_{p\mid N}\log N/\log p\ll(\log N)^2$; for each, the pairs
$(m,k)$ with $mk=N-n$ number $\le\tau(N-n)\ll N^{o(1)}$, and each term
is $\ll\log N$.
\end{proof}

By Lemma~\ref{lem:degen} we may drop the condition $(k,N)=1$ from
\eqref{eq:R1} at a cost of $O(N^{o(1)})$. This is essential: expanding
$\mathbf 1_{(k,N)=1}$ by M\"obius over $e\mid N$ would multiply the
number of Bombieri--Vinogradov calls by $2^{\omega(N)}=N^{o(1)}$ and
destroy the budget.

\subsection{Step 4: unfolding $\mu^2$ and the coprimality}

Insert $\mu^2(k)=\sum_{d^2\mid k}\mu(d)$ and
$\mathbf 1_{(k,m)=1}=\sum_{e\mid(k,m)}\mu(e)$ and truncate at
$D_0,E_0$ of \eqref{eq:trunc}.

\emph{Tail $d>D_0$.} Such terms have $d^2\mid k$, hence
$n=N-mk\equiv N\pmod{md^2}$; the number of $n\le N$ in that
progression is $\ll N/(md^2)+1$ and each $\Lambda\ll\log N$. Summing,
\[
  \ll\log N\sum_{m<M}\sum_{d>D_0}\Bigl(\frac{N}{md^2}+1\Bigr)
  \ll\frac{N(\log N)^2}{D_0}+\sqrt{NM}
  \ll\frac{N}{(\log N)^{A}}+N^{1-\theta'/2}.
\]

\emph{Tail $e>E_0$.} Such terms have $e\mid k$ and $e\mid m$, hence
$n\equiv N\pmod{me}$, and
\[
  \ll\log N\sum_{m<M}\sum_{\substack{e\mid m\\ e>E_0}}
      \Bigl(\frac{N}{me}+1\Bigr)
  \ll\frac{N\log N}{E_0}\sum_{m<M}\frac{\tau(m)}{m}+M
  \ll\frac{N(\log N)^{3}}{E_0}+M .
\]

Both are $\ll N(\log N)^{-A}$. In the remaining range put
$L:=\lcm(d^2,e)$ and $q:=mL$; the conditions $d^2\mid k$, $e\mid k$
and $n=N-mk$ give $n\equiv N\pmod q$, and $mk<N$,
$mk\ge\max(N-t,mK)$ give $1\le n\le T_m(t)$ where
\begin{equation}\label{eq:Tm}
  T_m(t):=\min\bigl(t,\ N-mK\bigr).
\end{equation}
Therefore, with $\psi(y;q,a)=\sum_{n\le y,\,n\equiv a\,(q)}\Lambda(n)$,
\begin{equation}\label{eq:R2}
  R(t)=\sum_{m<M}\mu(m)\sum_{d\le D_0}\mu(d)
         \sum_{\substack{e\mid m\\ e\le E_0}}\mu(e)\,
         \psi\bigl(T_m(t);\,q,\,N\bigr)
         \;+\;O\!\left(\frac{N}{(\log N)^{A}}\right).
\end{equation}
By Lemma~\ref{lem:degen} we may further restrict to $(q,N)=1$: if
$g=(q,N)>1$ then $g\mid n$ forces $n$ to be a power of a prime
dividing $N$, and the total over the $\ll MD_0E_0\tau$-many triples is
$\ll N^{1-\theta'+o(1)}$.

\begin{note}[the false-positive trap]\label{note:trap}
Restricting the \emph{main terms} to classes with $(q,N)=1$ is not
cosmetic. Assigning a main term $T_m/\varphi(q)$ to a degenerate class
shifts the density of the whole computation by the factor
$N/\varphi(N)$. Carried into the $\log k$ branch of
Section~\ref{sec:C}, that error produces an apparent \emph{refutation}
of $EH_\mu$. It is recorded here because it is the most plausible
false positive in this circle of ideas.
\end{note}

Note the size of the modulus: $q=m\lcm(d^2,e)\le MD_0^2E_0$, so by
\eqref{eq:trunc}
\begin{equation}\label{eq:level}
  q\;\le\;N^{1-\theta'}(\log N)^{4A+8}\;=\;N^{1/2-\delta}(\log N)^{4A+8}
  \;\le\;N^{1/2-\delta/2}
\end{equation}
for $N$ large. This is the level at which Bombieri--Vinogradov is
applied.

\subsection{Step 5: Bombieri--Vinogradov}

\begin{lemma}[$\tau$-weighted Bombieri--Vinogradov]\label{lem:BV}
For all $A,B>0$ there is $C=C(A,B)$ such that for
$Q\le N^{1/2}(\log N)^{-C}$,
\[
  \sum_{q\le Q}\tau_3(q)^{B}\ \max_{y\le N}\ \max_{(a,q)=1}
  \Bigl|\psi(y;q,a)-\frac{y}{\varphi(q)}\Bigr|
  \;\ll_{A,B}\;\frac{N}{(\log N)^{A}}.
\]
\end{lemma}

\begin{proof}
Write $\mathcal E_q$ for the inner maximum. By Cauchy--Schwarz,
$\sum_q\tau_3^B\mathcal E_q\le
(\sum_q\tau_3^{2B}\mathcal E_q)^{1/2}(\sum_q\mathcal E_q)^{1/2}$. For
the first factor use the trivial bound
$\mathcal E_q\ll(N/\varphi(q))\log N$, giving $\ll N(\log N)^{C_1(B)}$;
for the second use Bombieri--Vinogradov \cite{MV} with exponent
$2A+C_1(B)$.
\end{proof}

A modulus $q$ arises in \eqref{eq:R2} from at most
$\sum_{m\mid q}\tau(q/m)^2\le\tau(q)^3\le\tau_3(q)^3$ triples
$(m,d,e)$: given $m\mid q$ there are at most $\tau(q/m)$ choices of
$d$ with $d^2\mid q/m$ and at most $\tau(q/m)$ of $e$ with
$\lcm(d^2,e)=q/m$. Writing $\psi(T_m;q,N)=T_m/\varphi(q)+\mathcal E$
and applying Lemma~\ref{lem:BV} with $B=3$ and \eqref{eq:level},
\begin{equation}\label{eq:R3}
  R(t)\;=\;\mathrm{MT}(t)\;+\;O\!\left(\frac{N}{(\log N)^{A}}\right),
  \qquad
  \mathrm{MT}(t)=\sum_{\substack{m<M\\ (m,N)=1}}
       \mu(m)\,T_m(t)\,c_{D_0,E_0}(m),
\end{equation}
where
$c_{D_0,E_0}(m)=\sum_{d\le D_0}\sum_{e\mid m,\,e\le E_0}
\mu(d)\mu(e)\mathbf 1_{(d,N)=1}/\varphi(m\lcm(d^2,e))$.

\subsection{Step 6: the density}

\begin{lemma}\label{lem:density}
Let $m$ be squarefree with $(m,N)=1$ and put
\[
  c(m):=\sum_{d\ge1}\ \sum_{e\mid m}
     \frac{\mu(d)\mu(e)\mathbf 1_{(d,N)=1}}{\varphi(m\lcm(d^2,e))}.
\]
Then $c(m)=\AAA(N)\lambda(m)/m$ with $\AAA(N)$ as in \eqref{eq:AN} and
$\lambda(m)=\prod_{p\mid m}\bigl(1-\tfrac{1}{p(p-1)}\bigr)^{-1}$.
Moreover $c_{D_0,E_0}(m)=c(m)+O(1/(mD_0))$.
\end{lemma}

\begin{proof}
The sum is multiplicative. If $p\nmid m$ then $p\nmid e$ and the local
factor is $1-1/\varphi(p^2)=1-1/(p(p-1))$ for $p\nmid N$, and $1$ for
$p\mid N$. If $p\mid m$ (so $p\nmid N$) the $p$-part of
$m\lcm(d^2,e)$ is $p^{1+\max(2v_p(d),v_p(e))}$, and the four choices
$(v_p(d),v_p(e))\in\{0,1\}^2$ contribute
\[
  \frac{1}{p-1}-\frac{1}{p(p-1)}-\frac{1}{p^2(p-1)}+\frac{1}{p^2(p-1)}
  \;=\;\frac{1}{p}.
\]
Hence $c(m)=\prod_{p\mid m}p^{-1}\cdot
\prod_{p\nmid m,\ p\nmid N}(1-\tfrac{1}{p(p-1)})
=m^{-1}\AAA(N)\lambda(m)$. The truncation error is the tail
$\sum_{d>D_0}1/\varphi(md^2)\ll1/(mD_0)$ and similarly for $e$.
\end{proof}

The load-bearing line is the local factor. Equivalently, for squarefree
$m$,
\[
  \sum_{g\mid m}\frac{\mu(g)}{\varphi(m/g)\,g\,\varphi(g)}=\frac1m :
\]
the local factor is exactly $p^{-1}$, so the density exponent is
exactly $1$ and $1/\zeta$ occurs to the \emph{first} power in
Lemma~\ref{lem:mu} below. A non-integral exponent would give, by
Selberg--Delange, only $(\log x)^{-c}$ for a fixed $c$ in
Lemma~\ref{lem:mu}, and Theorem~\ref{thm:A} would be false. The
identity is verified in exact rational arithmetic for all squarefree
$m<400$, with zero mismatches
(\texttt{audit\_density\_identity.py}).

\begin{lemma}\label{lem:mu}
Let $f(m):=\mu(m)\lambda(m)\mathbf 1_{(m,N)=1}$ and
$G(x):=\sum_{m\le x}f(m)/m$. Then
\[
  G(x)\;\ll\;\frac{N}{\varphi(N)}\,e^{-c\sqrt{\log x}}
   +x^{-1/4}e^{C\sqrt{\log N}}
   \qquad(x\ge2).
\]
\end{lemma}

\begin{proof}
Write $F(s)=\sum_m f(m)m^{-s}=\prod_{p\nmid N}(1-\lambda(p)p^{-s})$
and $F(s)=\zeta(s)^{-1}H(s)$, so $f=\mu*h$ with $h$ multiplicative,
$h(p^j)=1$ for $p\mid N$ and $h(p^j)=1-\lambda(p)=O(p^{-2})$ for
$p\nmid N$. Then
$G(x)=\sum_{b\le x}\frac{h(b)}{b}\sum_{a\le x/b}\frac{\mu(a)}{a}$. For
$b\le\sqrt x$ use $\sum_{a\le y}\mu(a)/a\ll e^{-c\sqrt{\log y}}$ (the
classical zero-free region) together with
$\sum_b|h(b)|/b\ll\prod_{p\mid N}(1-1/p)^{-1}\ll N/\varphi(N)$. For
$b>\sqrt x$ bound
$\sum_{b>y}|h(b)|/b\le y^{-1/2}\sum_b|h(b)|b^{-1/2}
\ll y^{-1/2}\prod_{p\mid N}(1-p^{-1/2})^{-1}\ll y^{-1/2}e^{C\sqrt{\log N}}$.
\end{proof}

Since $M=N^{1-\theta'}$ is a fixed power of $N$, the second term of
Lemma~\ref{lem:mu} is negligible at $x\asymp M$, and
$N/\varphi(N)\ll\log\log N$; so $G(x)\ll e^{-c'\sqrt{\log x}}$ in the
relevant range.

\subsection{Step 7: the main term dies, uniformly in $t$}

\begin{proposition}\label{prop:MT}
$\displaystyle\sup_{1\le t<N}|\mathrm{MT}(t)|\ll Ne^{-c\sqrt{\log N}}$.
\end{proposition}

\begin{proof}
By Lemma~\ref{lem:density},
$\mathrm{MT}(t)=\AAA(N)\sum_{m<M}\frac{f(m)}{m}T_m(t)
+O(M\log N/D_0)$ with $T_m(t)$ as in \eqref{eq:Tm}. Put
$m_0:=(N-t)/K$, so $T_m(t)=t$ for $m\le m_0$ and $T_m(t)=N-mK$ for
$m>m_0$; in particular $x\mapsto T_x(t)$ is non-increasing, is
constant on $[1,m_0]$, has derivative $-K$ on $(m_0,M)$, and
$T_M(t)=0$. Abel summation against $G$ of Lemma~\ref{lem:mu} gives
\[
  \sum_{m<M}\frac{f(m)}{m}T_m(t)
  = G(M)\,T_M(t)+K\!\int_{m_0}^{M}\!G(x)\,dx
  = K\!\int_{m_0}^{M}\!G(x)\,dx .
\]
Hence, uniformly in $t$,
\[
  \Bigl|\sum_{m<M}\frac{f(m)}{m}T_m(t)\Bigr|
  \le K\int_{1}^{M}|G(x)|\,dx
  \ll K\int_1^M e^{-c\sqrt{\log x}}dx
  \ll KMe^{-c'\sqrt{\log M}}\ll Ne^{-c''\sqrt{\log N}},
\]
using $\log M\asymp\log N$. The truncation error is
$\ll M\log N/D_0=o(N(\log N)^{-A})$.
\end{proof}

The cancellation is genuine and global: the term $m=1$ alone
contributes $T_1(t)\asymp N$, so the smallness of $\mathrm{MT}$ is
entirely due to the cancellation in $\sum f(m)/m$.

\subsection{Conclusion}

Combining \eqref{eq:split0}, \eqref{eq:meanterm}, \eqref{eq:P},
\eqref{eq:R3} and Proposition~\ref{prop:MT}:
\[
  T_1(t)=\underbrace{P(t)}_{N^{o(1)}}
         -\underbrace{\mathrm{MT}(t)}_{\ll Ne^{-c\sqrt{\log N}}}
         -\underbrace{O\!\left(N(\log N)^{-A}\right)}_{\text{BV}}
         -\underbrace{C(t)B(K)}_{\ll Ne^{-c\sqrt{\log N}}},
\]
uniformly in $t<N$. This proves Theorem~\ref{thm:A}, and exhibits
Bombieri--Vinogradov as the only ingredient that does not give an
exponential saving. The bound cannot be improved to
$Ne^{-c\sqrt{\log N}}$ by this argument: the Bombieri--Vinogradov
input yields $N(\log N)^{-A}$ for each fixed $A$ and no more, because
the Siegel--Walfisz range in its proof is $q\le(\log N)^{B}$ with $B$
fixed. \qed

% =====================================================================
\section{Proof of Corollary~\ref{cor:B}}

Fix $k$ and set
$a_n:=\Lambda(n)\mu(N-n)\bigl(\mathbf 1_{n\equiv N\,(k)}-1/\varphi(k)\bigr)$,
so that $\Emu(t;k)=\sum_{n\le t}a_n$ and the bracket in \eqref{eq:E4}
is $\sum_{n<N}a_n\log(N-n)$. Since $\log(N-n)$ vanishes at $n=N-1$ and
has derivative $-1/(N-t)$, Abel summation gives
\[
  \sum_{n<N}a_n\log(N-n)\;=\;\int_{1}^{N-1}\frac{\Emu(t;k)}{N-t}\,dt .
\]
The weight $\log(N-n)$ does not depend on $k$, so multiplying by
$\mu(k)$ and summing over $k<K$ with $(k,N)=1$ (a finite sum) gives
the exact identity
\[
  E_4(\alpha)\;=\;\int_{1}^{N-1}\frac{T_1(t)}{N-t}\,dt .
\]
Hence $|E_4(\alpha)|\le\bigl(\sup_{t<N}|T_1(t)|\bigr)\log N
\ll_A N(\log N)^{1-A}$ by Theorem~\ref{thm:A}. As $A$ is arbitrary
this is Corollary~\ref{cor:B}. Huang--Li's Lemma~4 is invoked only to
bound $E_4$; with the above it is not needed, and the only remaining
appearance of $EH_\mu$ in their \S3 is through $E_3(\alpha)$. \qed

% =====================================================================
\section{Proof of Theorem~\ref{thm:C}: the permanent closure}%
\label{sec:C}

Take the weight $w_k=\log k$, i.e. the functional $E_3$ of
\eqref{eq:E3}. Steps~0--1 are unchanged, and the complete divisor sum
of Lemma~\ref{lem:complete} is replaced by the following.

\begin{lemma}\label{lem:completelog}
For squarefree $u$, write $u=u_Nu'$ where $u_N$ is the largest divisor
of $u$ composed of primes dividing $N$. Then
$\sum_{k\mid u,\ (k,N)=1}\mu(k)\log k=-\Lambda(u')$.
\end{lemma}

\begin{proof}
The sum is over $k\mid u'$ and equals $-\Lambda(u')$ by
$\mu*\log=\Lambda$.
\end{proof}

Hence the complete piece of $E_3$ is
\[
  -\sum_{1\le u<N}\Lambda(N-u)\,\mu(u)\,\Lambda(u')
  \;=\;-\sum_{\substack{u<N,\ (u,N)=1}}\Lambda(N-u)\mu(u)\Lambda(u)
        +O(N^{o(1)}\log^2 N).
\]
Now $\mu(u)\Lambda(u)$ is supported on primes $u=p$, where it equals
$-\log p$; prime powers $u=p^\ell$, $\ell\ge2$, have $\mu(u)=0$. So
the complete piece equals
\begin{equation}\label{eq:goldbachback}
  \sum_{p<N}\Lambda(N-p)\log p\;+\;O\bigl(N^{o(1)}\log^2N\bigr)
  \;=\;\sum_{n<N}\Lambda(n)\Lambda(N-n)+O\bigl(N^{1/2+\varepsilon}\bigr),
\end{equation}
the binary Goldbach sum itself. The door is free: divisor switching
costs nothing, and immediately behind it stands the object one was
trying to avoid.

Two further terms have to be evaluated, and neither vanishes.

\medskip\noindent
\emph{(i) The residual main term.} Repeating Steps~3--6 with the extra
factor $\log k=\log((N-n)/m)$ produces the main term
$\AAA(N)\sum_{m<M}\mu(m)\lambda(m)m^{-1}\int_{mK}^{N}\log(v/m)\,dv$.
Substituting $v=mt$ evaluates the integral as
$N\log(N/m)-N-mK\log K+mK$, of which only the $-N\log m$ piece
survives the cancellation: by Lemma~\ref{lem:mu} the coefficients
$\sum f(m)/m$ and $\sum f(m)$ both die, and what is left is
$-\AAA(N)\,N\,\widetilde G(1)$ with
$\widetilde G(1)=\lim_x\sum_{m\le x}\mu(m)\lambda(m)
\mathbf 1_{(m,N)=1}\log m/m$. The constant is fixed by
\begin{equation}\label{eq:Gtilde}
  \AAA(N)\,\widetilde G(1)\;=\;-\SS(N),
\end{equation}
so that the residual main term is $+\SS(N)\,N+O(N(\log N)^{-A})$.

The sign in \eqref{eq:Gtilde} is forced, and it is the whole identity:
with $F(s)=\sum_m f(m)m^{-s}=\zeta(s)^{-1}H(s)$ as in
Lemma~\ref{lem:mu} one has $\widetilde G(1)=-F'(1)=-H(1)<0$, while
$\AAA(N)H(1)=\SS(N)$ identically --- the local factors pair as
\[
  \Bigl(1-\tfrac{1}{p(p-1)}\Bigr)
  \Bigl(1-\tfrac{1}{(p^2-p-1)(p-1)}\Bigr)
  =1-\tfrac{1}{(p-1)^2}
\]
at $p\nmid N$, and as $p/(p-1)$ at $p\mid N$. Numerically,
$\AAA(N)\widetilde G(x)=-1.760250$ at $x=4\cdot10^6$ against
$\SS(N)=1.760432$ (\texttt{audit\_E3\_constant.py}).

\medskip\noindent
\emph{(ii) The subtracted mean term.} It carries the factor
\[
  B_{\log}(K):=\sum_{\substack{k<K\\ (k,N)=1}}
     \frac{\mu(k)\log k}{\varphi(k)}
  \;=\;-\SS(N)+O\bigl(e^{-c\sqrt{\log K}}\log K\bigr),
\]
which is exactly Huang--Li's Lemma~1: subtracting
$\sum_{k\le K}\mu(k)\varphi(k)^{-1}\log(K/k)
 =\SS(N)+O(e^{-c_1\sqrt{\log K}})$ from
$\log K\cdot\sum_{k\le K}\mu(k)/\varphi(k)
 =O(\log K\,e^{-c_1\sqrt{\log K}})$ gives the claim. Independently:
with $f(s)=\prod_{p\nmid N}\bigl(1-\frac{p^{-s}}{p-1}\bigr)$ one has
$f(s)=\zeta(s+1)^{-1}h(s)$ with $h(0)=\SS(N)$, so
$B_{\log}(\infty)=-f'(0)=-\SS(N)$. Since the mean term is
$-C(N)B_{\log}(K)$ with $C(N)=\sum_{n<N}\Lambda(n)\mu(N-n)$, it
contributes $+\SS(N)\,C(N)$.

\medskip
Collecting \eqref{eq:goldbachback}, (i) and (ii):
\[
  E_3(\alpha)=\sum_{n<N}\Lambda(n)\Lambda(N-n)
    -\SS(N)N+\SS(N)C(N)+O_A\!\left(\frac{N}{(\log N)^{A}}\right),
\]
which is the assertion of Theorem~\ref{thm:C}. Since $\SS(N)\asymp1$,
the bound $E_3(\alpha)\ll_A N(\log N)^{-A}$ is therefore equivalent to
\[
  \tilde r(N)=\SS(N)\bigl(N-C(N)\bigr)+O_A\!\left(N(\log N)^{-A}\right),
\]
which is Huang--Li's equation~(22). Given that, $\tilde r(N)$ is
positive for large even $N$ --- binary Goldbach --- because
$|C(N)|\le\AAA(N)N(1+o(1))$ with $\AAA(N)<1$; and
$\tilde r(N)\sim\SS(N)N$ holds if and only if $C(N)=o(N)$, which is
not implied. \qed

\begin{note}[the identity is not accessible to computation]%
\label{note:thetasweep}
The finite-$N$ residual
$R(N,\theta')=\bigl|E_3(N;\theta')-(\tilde r(N)-\SS(N)(N-C(N)))\bigr|/N$
does not become small in any accessible range, and it cannot be made
small by either free parameter. Over
$N=2\cdot10^5,\,4\cdot10^5,\,8\cdot10^5$ at $\theta'=0.56$ it reads
$0.4558,\,0.3729,\,0.3108$, while the right-hand side it is compared
against is $\asymp10^{-3}N$. Swept over $\theta'\in[0.51,0.95]$, $R$
is essentially monotone \emph{increasing} --- at $N=8\cdot10^5$ from
$0.170167$ at $\theta'=0.51$ to $4.991178$ at $\theta'=0.90$ ---
because the finite-$N$ error is dominated by the main-term
cancellation over $m<M=N^{1-\theta'}$, which $\theta'\to1$ destroys.
The best $\theta'$ is the smallest admissible one, just above $1/2$.
Increasing $N$ points the wrong way as well: at the best $\theta'$ the
ratio of error to signal reads $15.19,\,18.38,\,26.84$ at the three
$N$ above, because $C(N)=o(N)$ being true kills the right-hand side
fast while the finite-$N$ error dies only at a slow power of $\log$.
So Theorem~\ref{thm:C} is a statement no computation confirms, at any
$N$; the proof is the only access to it
(\texttt{lab\_theta\_sweep.py}).
\end{note}

\begin{note}
The mechanism is structural, not accidental. The identity
$\mu*\log=\Lambda$ is the starting point of the Huang--Li argument
(their~(10)); it is therefore also what any divisor switch inside
their $\log k$-weighted functional must return. No choice of
$\theta'$, truncation, or smoothing can circumvent an identity. What
the identity does not fix is the \emph{strength} at which $E_3$ must
be bounded, and that is Section~\ref{sec:strength}.
\end{note}

% =====================================================================
\section{How strong the surviving demand really is}\label{sec:strength}

\begin{proof}[Proof of Proposition~\ref{prop:onesided}]
By Theorem~\ref{thm:C},
$\tilde r(N)=\SS(N)(N-C(N))+E_3(\alpha)+O_A(N(\log N)^{-A})$. The
triangle inequality gives
$|C(N)|\le U(N):=\sum_{n<N}\Lambda(n)\mu^2(N-n)$, and
$U(N)=\AAA(N)N(1+o(1))$ by Mirsky's theorem on squarefree shifted
primes \cite{Mir49} together with partial summation. Hence
$\SS(N)(N-C(N))\ge\SS(N)(1-\AAA(N))N(1+o(1))$, and $\AAA(N)<1$ for
every $N$, so the right-hand side is positive. Substituting and using
that the $O_A$ term may be taken below any fixed power of $N/\log N$
gives $\tilde r(N)>0$ under \eqref{eq:onesided}.

For the size of the threshold, $\AAA(N)=\prod_{q\nmid N}
(1-\tfrac1{q(q-1)})$ is largest when $N$ carries the most small
primes, so $1-\AAA(N)$ is smallest at primorials, where
$1-\AAA(N)\asymp\sum_{q>p}\tfrac1{q(q-1)}\asymp1/(p\log p)$ with
$p\asymp\log N$; and $\SS(N)\gg1$.
\end{proof}

\begin{proof}[Proof of Proposition~\ref{prop:nolog}]
Immediate from Theorem~\ref{thm:C}, $|E_3(\alpha)|\le B(N)$, and
$|C(N)|\le\AAA(N)N(1+o(1))$ as above:
$\tilde r(N)\ge\SS(N)(1-\AAA(N))N(1+o(1))-B(N)
+O_A(N(\log N)^{-A})$.
\end{proof}

\begin{measurement}[the threshold, and the margin]\label{meas:margin}
Over even $N\le1.6\cdot10^7$ the threshold constant
$\SS(N)(1-\AAA(N))$ has median $0.333459$ and $0.1$-percentile
$0.119639$; its minimum is $0.060890$, attained at
$N=9\,699\,690=2\cdot3\cdots19$, the largest primorial in range, and
the ten smallest values are all primorial-like. Multiplied by
$\log N\log\log N$ the minimum over $N\ge10^5$ is $2.482019$, attained
at $N=510\,510$; over $N\ge10^3$ it is $1.737799$ at $N=2310$; over
$N\ge16$ it is $0.810202$ at $N=30$. Every argmin is a primorial,
which is the mechanism of Proposition~\ref{prop:onesided} made
visible; the constant $c$ there is therefore stated as a constant and
not as $1$. The step that produces the margin is exact to five
decimals: $U(N)/(\AAA(N)N)$ has mean $0.999996$ with standard
deviation $0.000170$ over the top octave. Checked against the truth,
no even $N\in[10^5,1.6\cdot10^7]$ has
$\tilde r(N)<\SS(N)(1-\AAA(N))N$; the slack between them has minimum
$3.7038$, median $4.2902$ and maximum $95.0889$, the maximum at
$N=9\,699\,690$ --- so the margin is thinnest exactly where the
Goldbach count is largest
(\texttt{lab\_onesided\_margin.py}).
\end{measurement}

\begin{measurement}[what the weakening relocates]\label{meas:relocate}
At $\theta'=0.56$ and $N=2\cdot10^5$ through $3.2\cdot10^6$, the
quantity $B(N)$ of \eqref{eq:Bsum} satisfies $B(N)/N=
0.8086,\,0.7395,\,0.7303,\,0.6547,\,0.5916$ against a threshold of
$0.3745N$, so the ratio $B(N)/(\SS(N)(1-\AAA(N))N)$ falls
$2.159,\,1.975,\,1.950,\,1.748,\,1.580$. The object
Proposition~\ref{prop:nolog} constrains is therefore already within a
factor $1.6$ of the threshold at the top of the accessible range, and
falling; but a constant-factor bound valid for \emph{all} large $N$ is
what is needed, and no computation supplies that
(\texttt{lab\_onesided\_demand.py}).
\end{measurement}

Proposition~\ref{prop:onesided} does not reopen the divisor switch.
The no-go of \cite{P2} costs a factor
$\exp(c_1\sqrt{\tfrac12\log N})$, which exceeds $\log N\log\log N$ as
it exceeds every power of $\log N$; so the weaker threshold is still
out of reach of that design space.

% =====================================================================
\section{The correction $\Delta$ to equation~(18)}\label{sec:delta}

We record a defect in the published paper, together with its repair.

Huang--Li define
$S_2(\alpha)=\sum_{n<N}\Lambda(n)\tilde\Lambda_\alpha(N-n)$ with
$\tilde\Lambda_\alpha(u)=\sum_{d\mid u,\ d>\alpha}\mu(d)\log(1/d)$,
and substitute $k=u/d$, so that the constraint $d>\alpha$ becomes
\begin{equation}\label{eq:constraint}
  k\;<\;\frac{N-n}{\alpha},
\end{equation}
an \emph{$n$-dependent} bound. Their displayed equation~(18) reads
\[
  S_2(\alpha)\;=\;\sum_{k<\frac{N-1}{\alpha}}\mu(k)
    \sum_{\substack{n<N\\ n\equiv N\ (k)}}
    \Lambda(n)\mu(N-n)\log\Bigl(\frac{k}{N-n}\Bigr),
\]
in which \eqref{eq:constraint} has been replaced by the $n$-free bound
$k<(N-1)/\alpha$. The two differ. Writing $N-n=mk$, the condition
\eqref{eq:constraint} is $m>\alpha$, so the right-hand side of~(18)
contains, in addition to $S_2(\alpha)$, exactly the terms with
$m\le\alpha$. On those terms
$\mu(k)\mu(N-n)=\mu(m)\mu^2(k)\mathbf 1_{(k,m)=1}$ and
$\log(k/(N-n))=-\log m$, so
\begin{equation}\label{eq:Delta}
  S_2(\alpha)\;=\;\text{RHS of (18)}\;+\;\Delta,
  \qquad
  \Delta\;=\;\sum_{2\le m\le\alpha}\mu(m)\log m
    \sum_{\substack{k<(N-1)/\alpha\\ (k,m)=1}}\mu^2(k)\,\Lambda(N-mk).
\end{equation}
(The term $m=1$ drops out since $\log1=0$.) The trivial bound is
$\Delta\ll N(\log N)^2$, which exceeds the target
$O(N(\log N)^{-A})$, so $\Delta$ is not negligible and needs its own
treatment.

It does close, by the machinery already used above and under
hypotheses Huang--Li already assume. Indeed $\Delta$ has exactly the
shape of the residual \eqref{eq:R1}: the M\"obius factor sits on the
\emph{short} variable $m\le\alpha$, and the long variable $k$ carries
only $\mu^2\ge0$. Expanding $\mu^2(k)$ and $\mathbf 1_{(k,m)=1}$ as in
Step~4 reduces $\Delta$ to sums of $\Lambda$ over progressions to
moduli $\ll\alpha(\log N)^{4A+8}$, plus a main term
$\AAA(N)\sum_{m\le\alpha}\mu(m)\lambda(m)(\log m)\,T_m/m$ which is
$O(Ne^{-c\sqrt{\log N}})$ by Lemma~\ref{lem:mu} and partial summation.
Hence:
\begin{itemize}
\item In the Corollary-1 regime, $\alpha=N^{1-\theta'}<N^{1/2}$ and
  Bombieri--Vinogradov closes $\Delta$ \emph{unconditionally}.
\item In the general Theorem-1 regime, $\alpha\asymp N^{\theta}$ and
  $\Delta$ closes under the hypothesis
  $EH(N^{\theta}(\log N)^{2A+8})$, which is assumed there.
\end{itemize}
So \cite{HL}'s Theorem~1 and their Corollary~1 stand as stated; what
is missing is the treatment of $\Delta$.

% =====================================================================
\section{Relation to the literature}\label{sec:lit}

This section is a placement, not a survey, and it is deliberately
conservative about what is offered here as new.

\medskip\noindent
\textbf{Theorem~\ref{thm:A}.} The mechanism is standard. After the
divisor switch the M\"obius on the long variable squares away and the
surviving M\"obius sits on a variable of length at most
$N^{1/2-\delta}$, which is the classical configuration in which
Bombieri--Vinogradov applies. Every ingredient
(Bombieri--Vinogradov \cite{MV}; the Goldston--Y{\i}ld{\i}r{\i}m
estimate \cite{GY} as used in \cite{HL}; the elementary density
identity of Lemma~\ref{lem:density}) is off the shelf. The statement
is offered as a \emph{lemma about the Huang--Li reduction}, not as a
contribution to the theory of Bombieri--Vinogradov; its content is
that this particular consumption of $EH_\mu$ is unnecessary, and what
difficulty it has lies in the bookkeeping --- in particular that the
condition $(k,N)=1$ must be discarded by the degeneracy argument
(Lemma~\ref{lem:degen}) rather than expanded over $e\mid N$, and that
main terms may be assigned only to classes with $(q,N)=1$, the second
of which shifts the density by exactly $N/\varphi(N)$ if missed
(Remark~\ref{note:trap}).

\medskip\noindent
\textbf{Theorem~\ref{thm:C}.} We are not aware of the identity in the
literature. Its ingredients are again classical; what it does is
locate, exactly, where the Goldbach content of the Huang--Li reduction
sits --- in the weight $\log k$, by $\mu*\log=\Lambda$.

\medskip\noindent
\textbf{Beyond the square-root barrier.} Since \cite{HL} appeared,
Lichtman \cite{Li23} has shown that the primes have level of
distribution $66/107\approx0.617$ using triply well-factorable
weights, and has used it to improve upper bounds for Goldbach
representations --- the first use of a level beyond the square-root
barrier on that problem. That work does not collide with the present
one: it concerns the primes rather than
$\sum_{n<N}\Lambda(n)\mu(N-n)$, and it yields upper bounds rather than
the asymptotic that \cite{HL} targets. It is nevertheless the state of
the art against which any claim of the form ``beyond $1/2$'' must now
be measured.

A level of distribution $3/5$ for the M\"obius function with triply
well-factorable weights belongs to the earlier part of that line of
work, \cite{Li20}; we cite it here as the object the question below
refers to rather than as a result we have checked against the
published version. It is \emph{not} in \cite{Li23}, which does not
treat the M\"obius function at all --- the $3/5$ appearing there is
Maynard's prior record for the \emph{primes}, which \cite{Li23}
improves to $66/107$. Since \cite{HL} require $EH_\mu(N^{\theta'})$
for a single $\theta'>1/2$, and since Theorem~\ref{thm:C} identifies
$w_k=\log k$ as the weight carrying the Goldbach content, the
following question is well posed and, so far as we know, open:

\begin{quote}
Does the weight that the Huang--Li consumption actually requires lie
in the triply well-factorable class for which level $3/5$ is known for
$\mu$?
\end{quote}

Two cautions attach to it. First, $EH_\mu$ as stated in
\eqref{eq:EHmu} carries $\max_a|\cdot|$, whereas well-factorable
results bound signed weighted sums $\sum_q\lambda_qE(x;q,a)$; the
comparison must therefore be made against what the argument
\emph{consumes}, which is a signed sum, and not against the
absolute-value form as written. Second, the no-go of \cite{P2} closes
extraction \emph{by divisor switching}; the technology of
\cite{Li23} is the Deshouillers--Iwaniec spectral large sieve, a
different mechanism, and the no-go says nothing about it.

% =====================================================================
\section{Summary}

\begin{itemize}
\item Theorem~\ref{thm:A}: the M\"obius-weighted, fixed-class
correlation sum with weight $w_k=1$ is $\ll_A N(\log N)^{-A}$,
unconditionally, for any level $N^{\theta'}$ with $\theta'>1/2$.

\item Corollary~\ref{cor:B}: hence Huang--Li's $E_4$ (their Lemma~4)
needs no hypothesis, and the whole $EH_\mu$ demand of their reduction
collapses to the scalar $E_3$.

\item Theorem~\ref{thm:C}: that scalar is unconditionally equivalent
to Huang--Li's equation~(22), hence already gives binary Goldbach for
large even $N$, and gives the asymptotic $\tilde r(N)\sim\SS(N)N$
exactly when $C(N)=o(N)$.

\item Propositions~\ref{prop:onesided} and~\ref{prop:nolog}: Goldbach
itself needs only a one-sided bound on $E_3$ at a threshold
$\asymp N$, equivalently a constant-factor bound on the object
$EH_\mu$ bounds --- weaker than the consumed bound by a factor
$(\log N)^A$ for every $A$.

\item Section~\ref{sec:delta}: equation~(18) of \cite{HL} omits an
$n$-dependent constraint; the missing term $\Delta$ is exhibited and
closes under hypotheses already assumed.

\item Section~\ref{sec:notclaimed}: the net progress toward Goldbach
is zero.
\end{itemize}

% =====================================================================
\begin{thebibliography}{99}

\bibitem{GY} D.~Goldston and C.~Y{\i}ld{\i}r{\i}m,
\emph{Higher correlations of divisor sums related to primes~I},
Integers \textbf{3} (2003), A5.

\bibitem{HL} J.-J.~Huang and H.~Li, \emph{On the connection between
the Goldbach conjecture and the Elliott--Halberstam conjecture},
arXiv:2005.03811v2 [math.NT], 2022.

\bibitem{Li20} J.~D.~Lichtman, \emph{Primes in arithmetic
progressions to large moduli, II: well-factorable estimates},
arXiv:2006.07088 [math.NT], 2020.

\bibitem{Li23} J.~D.~Lichtman (with an appendix by S.~Drappeau),
\emph{Primes in arithmetic progressions to large moduli, and Goldbach
beyond the square-root barrier}, arXiv:2309.08522 [math.NT], 2023.

\bibitem{Mir49} L.~Mirsky, \emph{The number of representations of an
integer as the sum of a prime and a $k$-th power free integer},
Amer. Math. Monthly \textbf{56} (1949), 17--19.

\bibitem{MV} H.~Montgomery and R.~Vaughan, \emph{Multiplicative number
theory~I: classical theory}, Cambridge University Press, 2007.

\bibitem{P2} \emph{No weight extracts $\sum_{n<N}\Lambda(n)\mu(N-n)$:
a no-go for the divisor-switch route, with the exact identities of its
design space}, companion paper.

\bibitem{Pan} C.-D.~Pan, \emph{A new attempt on Goldbach conjecture},
Chinese Ann. Math. \textbf{3} (1982), 555--560.

\end{thebibliography}

% =====================================================================
\appendix
\section{Reproducibility}

Every numerical figure printed above comes from one of the following
scripts; each is standalone, runs under Python with numpy on a laptop,
prints its own pass/fail criteria, and writes one result file.

\begin{center}
\begin{tabular}{lll}
\hline
Script & Result file & Supports\\
\hline
\texttt{audit\_switch\_identity.py} &
 \texttt{audit\_switch\_identity.txt} & \eqref{eq:switch}\\
\texttt{audit\_density\_identity.py} &
 \texttt{audit\_density\_identity.txt} & Lemma~\ref{lem:density}\\
\texttt{audit\_E3\_constant.py} &
 \texttt{audit\_E3\_constant.txt} & \eqref{eq:Gtilde}\\
\texttt{lab\_theta\_sweep.py} &
 \texttt{lab\_theta\_sweep.txt} & Remark~\ref{note:thetasweep}\\
\texttt{lab\_onesided\_margin.py} &
 \texttt{lab\_onesided\_margin.txt} & Measurement~\ref{meas:margin}\\
\texttt{lab\_onesided\_demand.py} &
 \texttt{lab\_onesided\_demand.txt} & Measurement~\ref{meas:relocate}\\
\hline
\end{tabular}
\end{center}

No computation is used as a premise of any proof. The computations
verify exact identities (\eqref{eq:switch}, Lemma~\ref{lem:density}),
fix the sign of a constant that the analysis determines
(\eqref{eq:Gtilde}), or measure a quantity whose size is stated only
over the range measured.

\end{document}
